% Keep styling.sty beside this file. Compile with:
% latexmk -xelatex -synctex=1 template.tex (or, 
% for my purposes, compile to preview with VimtexCompile in neovim)
% if youre using VSC with some built in plugin you'll probably need to add 
% {
%   "latex-workshop.latex.recipe.default": "latexmk (xelatex)",
%   "latex-workshop.latex.autoBuild.run": "onSave",
%   "latex-workshop.view.pdf.viewer": "tab",
%   "latex-workshop.view.pdf.invertMode.enabled": "never"
% }
% also it is necessary to download the right fonts. Just look 
% at error messages to find which ones or something

\documentclass[11pt]{article}
\usepackage{styling} % Dark by default, with the clickable THEME switch.
% \usepackage{graphicx} % Uncomment if adding external images.

\title{Assignment six}
\author{Thobias K. H\o ivik | Remote Group}
\date{For October 20, 2026} % Use \date{} to hide the date.

\begin{document}
\maketitle
Yippe, new preamble. 

\begin{prob*}[Munkres 30.5(a)]
    Show that every metrizable space with a countable 
    dense subset has a countable basis. 
\end{prob*}

\begin{lem*}
    Every separable metric space is 
    second-countable. 
\end{lem*}
\begin{proof}
    Let $(Y,d)$ be a separable metric space. 

    Meaning there exists $D \subseteq Y$, countable, 
    for which 
    \[
        \overline D = Y. 
    \]
    Recall that $\overline D$ is the set of all 
    points $x \in Y$ such that every open 
    neighborhood of $x$ intersects 
    $D$: 
    \[
        U \cap D \neq \emptyset   
    \]
    for every open neighborhood $U$ of $x$. 

    We have 
    \[
        D = \{x_1,x_2,\ldots\}. 
    \]
    Take 
    \[
        \mathcal B = \{B_{q}(x_i) \bigm| i = 1,2,\ldots, \ 
        q \in \mathbb Q_{>0}\}. 
    \]
    These are open sets and the 
    set is countable since we are indexing over 
    $\mathbb N \times \mathbb Q_{>0}$ which is 
    countable. 

    It is not hard to see that this covers the space. 

    Suppose $x \in Y$. 
    Fix $x_i \in D$ and choose $q \in \mathbb Q_{>0}$ such that 
    \[
        q > d(x,x_i). 
    \]
    Then 
    \[
        x \in B_{q}(x_i). 
    \]

    Now suppose $V$ is an open set. 

    Then, for $x \in V$, there exists $\varepsilon > 0$ 
    so that 
    \[
        B_{\varepsilon}(x) \subseteq V. 
    \]

    Since $D$ is dense, there exists 
    $x_i \in D$ such that 
    \[
        d(x,x_i) < \varepsilon/3. 
    \]
    Choose $r \in \mathbb Q_{>0}$ such that 
    \[
        d(x,x_i) < r < \varepsilon - d(x,x_i). 
    \]
    Such a rational exists since 
    \[
        2d(x,x_i) < \frac{2\varepsilon}{3} < \varepsilon. 
    \]
    Then $x \in B_r(x_i)$, and for every $y \in B_r(x_i)$ 
    the triangle inequality gives 
    \[
        d(x,y) \leq d(x,x_i) + d(x_i,y) < d(x,x_i) + r < 
        \varepsilon. 
    \]
    Thus 
    \[
        x \in B_r(x_i) \subseteq B_{\varepsilon}(x) 
        \subseteq V. 
    \]

    So $\mathcal B$ is a countable basis for $Y$, making 
    $Y$ second-countable. 

\end{proof}

\begin{proof}[Proof of 30.5(a)]
    Suppose $X$ is a metrizable space with a countable dense 
    subset. That is, 
    \[
        X \cong Y
    \]
    for some metric space $(Y,d)$ and there exists 
    $D\subseteq X$ such that 
    \[
        \overline D = X
    \]
    with $|D| \leq \aleph_0$. 

    Let $\varphi : X \to Y$ denote 
    the homeomorphism between them. 
    Then 
    \begin{align*}
        \varphi(\overline D) &= \overline{\varphi(D)} \\ 
        \varphi(X) &= \overline{\varphi(D)} \\ 
        Y &= \overline{\varphi(D)} 
    \end{align*}
    so $\varphi(D)$ 
    is dense in $Y$ and, since $\varphi$ is bijective, 
    $|\varphi(D)| \leq \aleph_0$, so 
    $\varphi(D)$ is a countable dense subset of $Y$. 

    Hence $Y$ is separable.

    By the lemma above, 
    $Y$ is second-countable. 

    Thus there exists a countable basis 
    \[
        \mathcal B_Y = \{B_1,B_2,\ldots\}
    \]
    of $Y$. 

    Then, for any open set $U$ in 
    $X$, since $X \cong Y$, we have 
    \[
        U = \varphi^{-1}(V)
    \]
    for open $V\subseteq Y$. 
    But 
    \[
        V = \bigcup_{i\in I} B_i 
    \]
    for some $I \subseteq \mathbb N$ 
    so 
    \begin{align*}
        U &= \varphi^{-1}\left(\bigcup_{i=1}^\infty B_i\right) \\ 
          &= \bigcup_{i=1}^\infty \varphi^{-1}(B_i). 
    \end{align*}
    Thus 
    \[
        \mathcal B = \{\varphi^{-1}(B_1), \varphi^{-1}(B_2), 
        \ldots\}
    \]
    is a countable basis of $X$, hence 
    $X$ is second-countable. 

\end{proof}


\end{document}
